\section{Recovering the radial matrix from a scalar determinant}
\label{cm-section}

The full matrix restriction can be deduced from the scalar Capelli
restriction.  Three commutators supply its other entries.  This avoids
three further expansions of bordered minors and, on the group itself,
introduces no division by a leading minor.

\subsection{The contracted matrix as commutators}

Use the quantum minors and signed complementary minors
$D(u),C_{ij}(u),N_{ilja}(u)$ of Section~\ref{rq-section}.
Let $n\geq2$, with generators satisfying
$[E_{ab},E_{cd}]=\delta_{bc}E_{ad}-\delta_{da}E_{cb}$.
Suppose that the commuting multiplication operators
$y_j,\eta_i$ satisfy
\[
 [E_{ab},y_j]=-\delta_{aj}y_b,
 \qquad [E_{ab},\eta_i]=\delta_{bi}\eta_a.
\]
For constant vectors $x=(x_j)$ and $\xi=(\xi_i)$ put
\[
 Z=\sum_j x_j\eta_j,\qquad B=\sum_i\xi_i y_i,
 \qquad [A,B]=AB-BA.
\]
Here $B$ denotes a multiplication operator; the letter does not
assert complex conjugacy.  In the unitary application below, $B=\bar Z$.
Define
\[
 \mathcal T(u)=
 \begin{pmatrix}
 D(u)&-H_\eta(u-1;x)\\
 -H_y(u;\xi)&-H_2(u-1;x,\xi)
 \end{pmatrix},
\]
with the contractions defined in Section~\ref{rq-section}.

\begin{theorem}\label{cm-commutator-matrix}
With multiplication coefficients placed to the left of the minors,
\begin{equation}\label{cm-four-commutators}
 \mathcal T(u)=
 \begin{pmatrix}
 D(u)&-[D(u),Z]\\
 {[D(u),B]}&-[[D(u),Z],B]
 \end{pmatrix}.
\end{equation}
This is an identity of differential operators on their common smooth
domain.  It requires neither a nonzero Gram determinant nor an inverse
spectral polynomial.
\end{theorem}

\begin{proof}
The determinant--vector identity \eqref{rq-det-vector} gives
\[
 [D(u),B]=-
       \sum_{i,j}\xi_i y_jC_{ij}(u)=-H_y(u;\xi).
\]
For the other first commutator, expand the determinant in column $j$,
putting its column factor first.  Every factor in the complementary
minor commutes with $\eta_j$, since its column differs from $j$.
Its commutator is therefore
\[
 [D(u),\eta_j]=\sum_i\eta_i C_{ij}(u-1).
\]
The shift $u-1$ is the parameter of the complementary minor after
the first column has been removed.  Multiplication by $x_j$ and
summation prove $[D(u),Z]=H_\eta(u-1;x)$.

The analogous row expansion of a complementary minor gives
\begin{equation}\label{cm-cofactor-vector}
 [C_{ij}(v),y_l]=-
                \sum_a y_aN_{ilja}(v).
\end{equation}
For $l=i$ both sides vanish.  Otherwise place row $l$ first
in the row expression for that minor.  The other rows commute
with $y_l$, and the remaining commutator is $-y_a$.
The two omission signs are precisely those in $N_{ilja}$.
The parameter of the remaining minor is $v$: the row expression
used here places its first row factor at the lowest parameter,
as in the proof of \eqref{rq-det-vector}.

All the multiplication coefficients commute.  Hence
\begin{align*}
 [[D(u),Z],B]
 &=-\sum_{i,j,l,a}\eta_i x_j\xi_l y_aN_{ilja}(u-1)\\
 &=\sum_{i,l,a,j}\eta_i\xi_l y_a x_jN_{ilaj}(u-1)
   =H_2(u-1;x,\xi).
\end{align*}
The second equality exchanges the two omitted column indices.
To see the last equality, group the four summands for each
unordered row pair and each unordered column pair.  Their factors
are $(\eta_i\xi_l-\eta_l\xi_i)(y_a x_j-y_jx_a)$, with exactly
the signed complementary minor in the definition of $H_2$.
Thus there is no extra factor of two or four.
\end{proof}

\begin{remark}\label{cm-small-dimension}
For $n=1$, take $C_{11}=1$ and $H_2=0$.  The identities become
$[u+E_{11},Z]=Z$, $[u+E_{11},B]=-B$, and
$[[u+E_{11},Z],B]=0$.  Thus the same formula holds with
$\mathcal T(u)=\bigl(\begin{smallmatrix}u+E_{11}&-Z\\-B&0
\end{smallmatrix}\bigr)$.
\end{remark}

\subsection{The full-group operators}

Let $(V,h)$ be any nondegenerate complex Hermitian space of dimension
$n\geq1$, with $h$ linear in its first argument, and let $x,z\in V$
be fixed.
For $g\in U(V,h)$ put
\[
 Z(g)=h(gx,z),\qquad
 y(g)=g^{-1}z,\qquad
 \eta_i(g)=h(gw_i,z),\qquad \xi_i=h(w_i,x).
\]
Then $B(g)=\sum_i\xi_i y_i(g)=h(g^{-1}z,x)=\overline{Z(g)}$.
The complexified right derivatives satisfy the hypotheses of
Theorem~\ref{cm-commutator-matrix}.  No definiteness assumption
on $V$, and no assumption on the signs of $h(x,x),h(z,z)$, is needed.

\begin{corollary}\label{cm-group-preservation}
The four entries of $\mathcal T(u)$ are globally smooth differential
operators with polynomial coefficients in the matrix entries of $g$
and their conjugates.  They commute with left and right translations
by subgroups fixing $z$ and $x$, respectively.  Their formal adjoints
have the same coefficient property.  After the real shift
$u=Y+(n-1)/2$, for real $Y$ their matrix of Haar formal adjoints
is their transpose.

Consequently these entries and all their finite compositions preserve
smooth compact support, a prescribed positive Gram cutoff, and the
spaces of rapidly decreasing smooth functions used in
Section~\ref{jf-section}.
\end{corollary}

\begin{proof}
The full Capelli determinant is central and is invariant under
conjugation, so its right differential operator commutes with all
left and right translations.  The coefficient $Z$ is invariant
under the translations in the statement, and so is $\bar Z$.
Formula~\eqref{cm-four-commutators} proves the claimed covariance.

The invariant derivatives have polynomial coefficients in matrix
coordinates.  Also $g^{-1}=H^{-1}g^*H$ on $U(V,h)$, where $H$
is a fixed Gram matrix.  Thus both $Z$ and $\bar Z$, and all
their invariant derivatives, have polynomial growth.  No inverse
coordinate or Gram factor occurs in \eqref{cm-four-commutators}.
The same is true of the Haar formal adjoints, because invariant
vector fields have zero divergence and adjoint formation differentiates
these same polynomial coefficients.

For the adjoint identity, $D(u)^\dagger=D(\bar u)$ follows from
the signed-minor adjoint rule used in Lemma~\ref{re-base-leading}.
Multiplication by $Z$ has adjoint multiplication by $\bar Z$.
Thus the two off-diagonal entries are adjoints.  The last entry
is self-adjoint, since $[Z,\bar Z]=0$ and the Jacobi identity give
$[[D,Z],\bar Z]=[[D,\bar Z],Z]$.
Finally the Leibniz rule proves preservation of rapid decay after
increasing finitely many derivative and decay indices.  Differential
operators do not enlarge support, so they also preserve the cutoff.
\end{proof}

\subsection{One scalar restriction gives the entire matrix restriction}

Let $c\in\mathbb C^\times$, let $p=c\partial_c$ and
$\bar p=\bar c\partial_{\bar c}$, and put
\[
 T(Y)=\begin{pmatrix}Y+p&-c\\c^{-1}&0\end{pmatrix},\qquad
 L(Y)=\begin{pmatrix}Y-\bar p&\bar c^{-1}\\-\bar c&0\end{pmatrix}.
\]
Let $t(Y)=\bigl(\begin{smallmatrix}a&b\\d&e\end{smallmatrix}\bigr)$
be a matrix of differential operators in independent variables.
Its entries may fail to commute with each other; they commute with
$c,\bar c,p,\bar p$.  Set $M=L(Y)t(Y)T(Y)$ and $C=M_{11}$.

\begin{lemma}\label{cm-scalar-sandwich}
One has
\begin{equation}\label{cm-sandwich-recovery}
 M=\begin{pmatrix}
 C&-[C,c]\\
 {[C,\bar c]}&-[[C,c],\bar c]
 \end{pmatrix}.
\end{equation}
\end{lemma}

\begin{proof}
Multiplication in the displayed order gives
\[
 C=(Y-\bar p)a(Y+p)+(Y-\bar p)b c^{-1}
                      +\bar c^{-1}d(Y+p)+\bar c^{-1}e c^{-1}.
\]
The only nonzero commutators with the coordinate multipliers
come from $[p,c]=c$ and $[\bar p,\bar c]=\bar c$.
Consequently
\begin{align*}
 -[C,c]&=-(Y-\bar p)ac-\bar c^{-1}dc=M_{12},\\
 [C,\bar c]&=-\bar c a(Y+p)-\bar c b c^{-1}=M_{21},\\
 -[[C,c],\bar c]&=\bar c a c=M_{22}.
\end{align*}
This proves all entries without interchanging operators belonging
to $t(Y)$.
\end{proof}

\begin{theorem}\label{cm-full-restriction}
Let $V=\mathbb C e_0\oplus W\oplus\mathbb C f_0$ be a Hermitian
space, where $e_0,f_0$ form a hyperbolic pair and $\dim W=d\geq2$.
Use the open-cell coordinates and character laws of
Section~\ref{rr-section}, with any fixed vectors $z,x\in W$.
Form the contracted matrix $t(Y)$ on $U(W)$ from these vectors,
with parameter $Y+(d-1)/2$.
On $U(V)$ form the contracted matrix $Q(Y)$ using the vector
$e_0$ on both sides and parameter $Y+(d+1)/2$.
For a smooth function $F$ with the two prescribed unipotent
character laws, write
\[
 (\mathcal RF)(c,m)=|c|^{d+1}F(w(c)m).
\]
Then every entry satisfies
\begin{equation}\label{cm-full-one-pair}
 \mathcal R(Q_{ab}(Y)F)
        =(L(Y)t(Y)T(Y))_{ab}\mathcal RF,
               \qquad a,b\in\{1,2\}.
\end{equation}
The identity is pointwise on the open cell.  It allows isotropic
$z$ and $x$ and arbitrary signature of $W$.
\end{theorem}

\begin{proof}
The scalar $(1,1)$ identity is the bordered determinant calculation
of Lemma~\ref{rr-radial-capelli}.  That calculation uses only the
nondegeneracy of $W$, its two character laws, the vanishing of their
central characters, and the diagonal derivative rules.  None of its
steps uses a sign for the norm of $z$ or $x$, or restricts the signature
of $W$.  More explicitly, with
$y=m^{-1}z$, $\eta_i=h(mw_i,z)$, and $\xi_i=h(w_i,x)$,
the character rules are
\[
 E_{e_0j}F=x_jF,\quad E_{if_0}F=\xi_iF,\quad
 (E_{ie_0}XF)(w(c)m)=c^{-1}\eta_i(XF)(w(c)m),
\]
\[
 (E_{f_0j}XF)(w(c)m)=\bar c^{-1}y_j(XF)(w(c)m).
\]
The first two are right character rules, the last two are left
character rules and hold for every right differential word $X$.
The central root terms vanish as in
\eqref{rc-mixed-right}--\eqref{rc-mixed-left}, and
$E_{e_0e_0}$ and $E_{f_0f_0}$ restrict to $p$ and $-\bar p$.
These statements use only the Hermitian pairing and the chosen
characters.  The bordered expansion therefore applies with full
parameter $u=Y+(d+1)/2$ and $k+1$ replaced by $d$.
The normalizing factor has
$p|c|^{d+1}=\bar p|c|^{d+1}=(d+1)|c|^{d+1}/2$.
Thus its conjugation gives exactly the two factors $Y+p$ and
$Y-\bar p$, and the middle parameter is $Y+(d-1)/2$.

On the larger group the matrix coefficient for $e_0,e_0$ is
$Z_0(g)=h(ge_0,e_0)$.  It is invariant under both integrating
unipotent groups, and $Z_0(w(c)m)=c$.  Hence $Z_0F$ and
$\bar Z_0F$ have the same character laws as $F$, and
\[
 \mathcal R(Z_0F)=c\mathcal RF,\qquad
 \mathcal R(\bar Z_0F)=\bar c\mathcal RF.
\]
Apply the scalar identity also to these products.  Subtraction
shows that $\mathcal R$ carries $[Q_{11},Z_0]$ to $[C,c]$,
$[Q_{11},\bar Z_0]$ to $[C,\bar c]$, and their mixed
commutator to $[[C,c],\bar c]$.  Theorem~\ref{cm-commutator-matrix}
and Lemma~\ref{cm-scalar-sandwich} now give all of
\eqref{cm-full-one-pair}.  The real scalar normalizing factor
commutes with both coordinate multipliers, so it introduces no
additional term in these commutators.
\end{proof}

\begin{corollary}\label{cm-chain-restriction}
Fix integers $k\geq1$ and $r\geq0$.  For a succession of $r$
hyperbolic pairs ending in a Hermitian space of dimension $k+1$,
choose the two distinguished vectors in each nonterminal remaining
space to be its first isotropic basis vector, and those in the
terminal space to be its two fixed vectors.  At each step use the
unipotent characters determined by the distinguished vectors of the
next remaining space.  These are the standard flag characters.
Let $t(Y)$ be the
terminal contracted matrix and let $T_j(Y),L_j(Y)$ act in the
$j$th independent nonzero coordinate.  The restriction of the full
initial matrix, with the normalizing factor $\gamma_r$ of
Section~\ref{jf-section}, namely
\[
 \gamma_r(c)=\prod_{j=0}^{r-1}|c_j|^{k+2(r-j)},
\]
is
\begin{equation}\label{cm-full-chain}
 L_0(Y)\cdots L_{r-1}(Y)\,t(Y)\,
                       T_{r-1}(Y)\cdots T_0(Y).
\end{equation}
This includes the empty chain.  All full-group entries and finite
words in them have the coefficient, support and adjoint properties
required for the Jacquet commutation argument of
Lemma~\ref{jf-pointwise-smooth}.
\end{corollary}

\begin{proof}
Apply Theorem~\ref{cm-full-restriction} one pair at a time.
At each nonterminal stage the distinguished vector is the first
isotropic vector of the remaining space.  The next stage's character
vectors may therefore be isotropic, as allowed by the theorem.
At step $j$ the larger space has dimension $k+1+2(r-j)$,
so its factor is $|c_j|^{k+2(r-j)}$.
These factors multiply to $\gamma_r$, and associativity
gives the indicated order of the matrices.  Operators in previously
separated variables commute with the next restriction.  For $r=0$
the formula is the terminal matrix itself.
The last assertion follows from Corollary~\ref{cm-group-preservation}.
Each entry commutes with the translations defining its Jacquet
equivariance, so every finite product does as well.
The dominated-differentiation and Haar-adjoint arguments already
proved in Sections~\ref{jp-section} and \ref{jf-section} therefore
apply to these operators, with the same positive Gram support.
\end{proof}

These results establish the full matrix restriction and its
compatibility with the existing Jacquet operators.  They give no
estimate uniform in the length of a chain.  The noncompact
residue-pairing application must still assemble its actual generated
vectors, convergence statements, and scalar weights; the global
orbital-energy and arithmetic divisibility estimates remain separate.
